%
% Template for Doctoral Theses at Uppsala 
% University. The template is based on    
% the layout and typography used for      
% dissertations in the Acta Universitatis 
% Upsaliensis series                      
% Ver 5.2 - 2012-08-08                  
% Latest version available at:            
%   http://ub.uu.se/thesistemplate            
%                                         
% Support: Wolmar Nyberg Akerstrom        
% Thesis Production           
% Uppsala University Library              
% avhandling@ub.uu.se                          
%                                         
%%%%%%%%%%%%%%%%%%%%%%%%%%%%%%%%%%%%%%%%%%%


\documentclass{stylefiles/UUThesisTemplate}

% Package to determine wether XeTeX is used
\usepackage{ifxetex}


\ifxetex
	% XeTeX specific packages and settings
	% Language, diacritics and hyphenation
	\usepackage[babelshorthands]{polyglossia}
	\setmainlanguage{english}
	\setotherlanguages{swedish}

	% Font settings
	\setmainfont{Times New Roman}
	\setromanfont{Times New Roman}
	\setsansfont{Arial}
	\setmonofont{Courier New}
\else
	% Plain LaTeX specific packages and settings
	% Language, diacritics and hyphenation
    % Use English and Swedish languages. 
	\usepackage[swedish,english]{babel} 

	% Font settings
	\usepackage{type1cm}
	\usepackage[utf8]{inputenc}
	\usepackage[T1]{fontenc}
	\usepackage{mathptmx}
	
	% Enable scaling of images on import
	\usepackage{graphicx}
\fi


% Tables
\usepackage{booktabs}
\usepackage{tabularx}
% For codes
\usepackage{caption}
\captionsetup[listing]{justification=centering, singlelinecheck=false}
\usepackage{amsmath,amssymb,mathtools}
\usepackage{minted}
\setminted{
	fontsize=\small,
	bgcolor=gray!3,
	frame=single,
	breaklines=true,
	autogobble=true
}

\usepackage{tcolorbox}       % boxed examples
\tcbset{colback=gray!6, colframe=gray!40, boxrule=0.4pt, arc=2pt}

% Document links and bookmarks
\usepackage{hyperref} 
\usepackage{cleveref}

\usepackage{csquotes} % For displayquote
\newcommand{\rfcdisplayquote}[1]{%
	\begingroup%
	\vskip .25em%
	\addtolength\leftmargini{-1em}%
	\begin{displayquote}%
		\textsl{\small #1}%
	\end{displayquote}%
	\vskip .15em%
	\endgroup%
	\noindent%
}

\newcommand{\rfc}[1]{\texttt{#1}}
\newcommand{\rfcm}[1]{\mathtt{\scriptstyle#1}}
% Numbering of headings down to the subsection level
\numberingdepth{subsection}

% Including headings down to the subsection level in contents
\contentsdepth{subsection}


% Uncomment to use a custom abstract dummy text
\abstractdummy{
	\begin{abstract}
		Testing network protocol implementations for conformance with their specifications is challenging due to protocol statefulness, complex dependencies across packet sequences, and large input spaces that hinder exhaustive exploration. 
		Subtle deviations from specifications can result in interoperability issues or security vulnerabilities, yet often escape detection by conventional testing techniques. 
		This thesis investigates how symbolic execution and related techniques can be adapted to address these challenges in a systematic and practical manner. 
		
		The first contribution of the thesis is a requirement-driven symbolic execution approach for testing stateful network protocol implementations against their specifications. 
		The approach guides symbolic execution toward execution paths that are relevant for violating protocol requirements expressed over sequences of packets. 
		By selectively making packet fields symbolic and constraining exploration to requirement-relevant paths, the approach enables the detection of violations that depend on subtle corner-case values.
		
		The second contribution is a monitor-based symbolic execution technique that externalizes requirement checking from the system under test. 
		Protocol requirements are encoded in lightweight monitors that observe packet exchanges, maintain interaction state, and guide symbolic exploration without intrusive modifications to protocol implementations. 
		This design improves modularity, supports reuse of requirements across implementations and versions, and enables early detection of requirement violations. 
		
		The third contribution explores differential symbolic execution as an oracle-free testing technique. 
		Instead of relying on explicitly encoded requirements, the approach compares observable behaviors across implementations or versions under identical symbolic inputs to identify discrepancies that may indicate nonconformance.
		
		The final contribution examines requirement-based testing of IoT protocol implementations using fuzzing and symbolic execution, with CoAP as a case study. 
		The study analyzes how both techniques can be adapted to check protocol requirements and compares their effectiveness and practical trade-offs in exposing specification violations. 
		
		The proposed techniques are evaluated on real-world implementations of widely deployed protocols, including DTLS, QUIC, and CoAP. 
		The evaluation demonstrates that the proposed techniques uncover numerous previously unknown bugs, specification violations, and interoperability issues.
	\end{abstract}
}


\begin{document}
\frontmatter
    % Creates the front matter (title page(s), abstract, list of papers)
    % for either a Comprehensive Summary or a Monograph.
    % Authors of Comprehensive Summaries use this front matter 
    \frontmatterCS 
    % Monograph authors use this front matter 
    %\frontmatterMonograph 
 
   % Optional dedication
   \dedication{Dedicated to my parents Hanieh \& Mohammad}
 
    % Environment used to create a list of papers
    \begin{listofpapers}
    	\item \textbf{Applying Symbolic Execution to Test Implementations of a Network Protocol Against its Specification} \label{ICST@2022} \\
    	Hooman Asadian, Paul Fiter\u{a}u-Bro\c{s}tean, Bengt Jonsson, and Konstantinos Sagonas. In \emph{IEEE Conference on Software Testing, Verification and Validation (ICST 2022)}. 70-81. \url{https://doi.org/10.1109/ICST53961.2022.00019}
    	
    	\item \textbf{Monitor-based Testing of Network Protocol Implementations Using Symbolic Execution}\label{ARES@2024} \\
    	Hooman Asadian, Paul Fiter\u{a}u-Bro\c{s}tean, Bengt Jonsson, and Konstantinos Sagonas. In \emph{Proceedings of the 19th International Conference on Availability, Reliability and Security ARES '24}. ACM, Article 17, 1-12.
    	\url{https://doi.org/10.1145/3664476.3664521}
    	
    	\item \textbf{Uncovering Flaws in Network Protocol Implementations Using Differential Symbolic Execution}\label{symdiff@under} \\
    	Hooman Asadian, Paul Fiter\u{a}u-Bro\c{s}tean, Bengt Jonsson, and Konstantinos Sagonas.
    	Draft
    	\url{}
    	
    	\item \textbf{Testing IoT Protocol Requirements Using Fuzzing and Symbolic Execution: Application to CoAP}\label{cscn@2024} \\
    	Hooman Asadian, Paul Fiter\u{a}u-Bro\c{s}tean, Bengt Jonsson, and Konstantinos Sagonas. In \emph{IEEE Conference on Standards for Communications and Networking (CSCN 2024)}. 48-54.
    	\url{https://doi.org/10.1109/CSCN63874.2024.10849709} 
    	
    \end{listofpapers}
    
    
    \begingroup
        % To adjust the indentation in your table of contents, uncomment and enter the widest numbers for each level
        %  E.g.  \settocnumwidth{widest chapter number}{widest section number}{widest subsection number}...{...}
       %  \settocnumwidth{5}{4}{5}{3}{3}{3}
        \tableofcontents
    \endgroup
    
    % Optional tables
    %\listoftables
    %\listoffigures

\mainmatter
    % This includes the "Instruction", "Problem and Solutions" and "Example" files. After reading it, remove it from Thesis.tex. 
    \input{chapters/introduction.tex}
    \input{chapters/protocol-testing.tex}
    \chapter{From Specifications to Oracles: Requirement-Driven Symbolic Execution}\label{sec:paper1}
Network protocols rely on precise specifications that describe how communicating parties should construct messages, how state is expected to evolve during an interaction, and how error conditions must be handled.
These specifications are expressed in RFC documents, where natural language is used to formulate requirements.
Many of these requirements involve relationships across multiple packets.
For example, a field in one message may depend on a value in a previous message, or a sequence must follow a specific ordering that is valid for the selected key exchange method.

Testing such requirements is challenging. 
Traditional fuzzers concentrate on triggering crashes and often operate without awareness of protocol state.
Standard symbolic execution (SE) tools are powerful for analyzing single inputs, yet they do not directly address multi packet interactions or stateful requirements.
When applied naively to protocol implementations, these tools either fail to explore relevant paths or become overwhelmed by the path explosion caused by large input structures.

The work presented in this chapter introduces a technique that applies SE in a requirement driven manner.
The aim is to systematically search for inputs that violate specific RFC requirements and to produce concrete packet sequences that expose the deviation in the implementation under test.
The objective is not limited to detecting crashes. 
It also focuses on revealing non-conformance flaws that may affect security or interoperability.
\section{Technique}
\subsection{Extracting Specification Requirements}
The first step is to extract protocol requirements from the RFC text and express them as logical predicates over the packets exchanged during a protocol session.
RFCs typically use particular keywords such as "MUST", "MUST NOT", "REQUIRED", "SHALL", and "SHALL NOT" to indicate the strictness of a requirement~\cite{RFC2119}.
These statements can be formulated into predicates that describe how packets must be structured and how they relate to one another.

A simple but representative example appears in the DTLS RFC~\cite[p. 13]{RFC6347}, which specifies the uniqueness of record sequence numbers within a session.
Each DTLS record contains a 48-bit \rfc{sequence\_number} field that uniquely identifies the record within a session.
This field is incremented for every record sent, usually starting from zero. 
The receiver uses sequence numbers to process records in the correct order and to detect duplicates, since any two records with the same \rfc{sequence\_number} represent a replay.
%
\rfcdisplayquote{%
For each received record, the receiver MUST verify that the record contains a sequence number that does not duplicate the sequence number of any other record received during the life of this session.}%
%
This requirement serves as a protection against replayed records (packets).
For a set of records $R$ received during a DTLS session, the requirement can be expressed in predicate logic as:
%
\begin{equation}\label{const:rsnuniqueness}\normalsize
	\forall r, r' \in R :\, r \neq r' \implies 
	r\rfcm{.sequence\_number}\, \neq\, r'\rfcm{.sequence\_number}
\end{equation}
%
This formula captures the essential rule that no two records may share the same sequence number.
Requirements of this kind, which relate fields across multiple packets, appear frequently in network protocol specifications. 
Once formulated, they provide a precise basis for guiding the SE engine toward the regions of the input space where potential violations may occur.

\subsection{Symbolic Execution}
Symbolic execution is applied to each tested implementation through a dedicated test harness that replays a pre-captured sequence of records from a valid DTLS handshake and passes them one at a time to the implementation under test. 
Each record is parsed into a structured representation, which enables SE to access and manipulate selected fields. 
The harness provides a deterministic environment by controlling the order of record processing and by removing nondeterministic behavior such as cookie verification or timeouts.

With the test harness in place, the next step is to select which fields should be made symbolic for checking a requirement.
In our technique, only those fields in a record that are relevant to the requirement being tested are made symbolic, while the remaining fields retain their concrete values from the pre-captured handshake.
This selective approach focuses SE on the parts of the input space that influence the requirement, while avoiding possible path explosion if entire records were made symbolic.

To ensure that SE concentrates on violating cases rather than on the vast space of valid inputs, the harness introduces assumptions that encode the negation of the requirement.
These assumptions restrict SE to inputs under which the requirement may be broken.
The final step is to check whether the implementation uses the invalid input in a forbidden way.
This can be done by checking whether protocol state progresses after reception of invalid records.
We approximate this by successful completion of the protocol handshake. 
A failing assertion is inserted at the point in the implementation code where the handshake finishes successfully.
If SE reaches this assertion, it means that the implementation completed the handshake even though it received an invalid record, which indicates a violation of the requirement.

The predicate in \Cref{const:rsnuniqueness} formulates the DTLS requirement that no two records in a session may share the same sequence number.
To apply our technique to evaluate this requirement for a server, the harness loads the sequence of pre-captured records that a DTLS server would receive during its handshake and marks the client-side \rfc{sequence\_number} fields as symbolic.
All other fields remain concrete.

An assume statement is then introduced that indicates that at least two symbolic sequence numbers are equal.
This directs SE to explore precisely those assignments that violate uniqueness.
To evaluate the effect of this invalid input, an assertion is placed at the code location that indicates successful completion of the handshake. 
If SE reaches this assertion, it means that the implementation completed the handshake even though it received a record with a duplicate sequence number. 
The assertion then fails, and the SE engine reports the concrete assignments that triggered this behavior.

\subsection{Test Case Construction and Validation}
When SE discovers a violating execution path, the engine provides concrete values for all symbolic fields involved.
To construct a test case, these values are substituted back into the structured record representations used by the test harness, while all remaining fields retain their original concrete values from the pre-captured handshake.
The resulting records are then serialized into byte arrays that follow the exact wire format of the protocol.

For the sequence number requirement, this means that SE produces a concrete assignment in which two records share the same sequence number. 
These values replace the symbolic fields in the corresponding records, creating a concrete record sequence that instantiates the violation of the predicate in \Cref{const:rsnuniqueness}.

These concrete record sequences are replayed against the unmodified implementation outside the SE environment.
Validation is done by observing whether the unmodified implementation accepts the sequence and completes the handshake.
Since handshake completion is used as an approximation for protocol progress, a successful handshake in the presence of an invalid record constitutes a confirmed violation of the corresponding requirement.
In the sequence number example, if the server completes its handshake despite the duplicate sequence number, the requirement is violated and the test case is confirmed.

\section{Summary of Results}
The technique was applied to four implementations of the DTLS protocol. These included OpenSSL, MbedTLS, and two variants of TinyDTLS. A collection of sixteen requirements from the DTLS RFC~\cite{RFC6347} were formulated and tested. These requirements ranged from record layer constraints, such as content type, version fields, and epoch handling, to handshake layer rules related to message ordering, fragmentation, and message sequence numbers.

The evaluation revealed a total of thirty six unique bugs, seven of which were classified as vulnerabilities. Several issues were memory related, including out of bounds accesses during message reassembly. Others concerned logical violations of the specification, such as incorrect validation of version fields, improper handling of repeated handshake messages, and failures to respect message ordering rules. Many issues were identified within a short symbolic exploration time, and all of them were acknowledged and fixed by the developers.

In addition to uncovering new issues, our technique also reproduced a known exploitable vulnerability in OpenSSL's function for reassembling fragmented messages.
This reproduction confirms the power and promise of our technique.

For a detailed description of the technique's implementation, the full set of requirements we evaluated, and a comprehensive analysis of the issues we discovered, the reader is referred to \textcolor{red}{Paper 1} included later in this thesis.
\section{Limitations}
The technique presented in this chapter has two main limitations that motivate the next contribution of this thesis.
The first limitation is that the code responsible for checking requirements must be directly inserted into the source code of the implementation under test.
This involves adding calls to a shared library at specific locations in the implementation.
Since these locations differ between implementations and between versions, the process requires a detailed understanding of the code base and must be repeated whenever the implementation changes.
This reduces the portability of the technique and increases the manual effort required to apply it.

The second limitation concerns the reliance on completing a protocol interaction.
In the case of DTLS, this means that SE must reach either the successful completion of the handshake or a definite failure triggered by the invalid input.
As a result, SE may be forced to explore many paths that lie far beyond the part of the execution that is relevant to the requirement being tested.
This unnecessarily increases the symbolic search space and can significantly slow down the detection of violating inputs.

These two limitations motivate the need for a technique that avoids direct instrumentation of the implementation and that can detect requirement violations without relying on the completion of a full protocol exchange.
The next chapter introduces a monitor based approach that addresses both issues by moving requirement checking outside the implementation and by enabling earlier detection of violation points.
    \chapter{Monitor-based Symbolic Execution for Stateful Protocols}\label{sec:paper2}
\section{Technique}
\subsection{Creating a Test Harness}
\subsection{Encoding Monitors as Programs}
\subsection{Invoking Monitors}
\subsection{Developing Protocol-specific Infrastructure}
\subsection{Applying Small Modifications to the SUT}
\section{Summary of Results}
\section{Limitations}
    \chapter{Beyond Oracles: Differential Symbolic Execution}\label{chapter:paper3}
The techniques presented in the previous chapters rely on the existence of an explicit oracle that characterizes correct protocol behavior.
In \Cref{chapter:paper1}, correctness is defined through logical predicates extracted from protocol specifications and checked by instrumenting the implementation under test.
In \Cref{chapter:paper2}, this notion is preserved, but the responsibility for checking requirements is moved outside the implementation and delegated to external monitors.
In both cases, symbolic execution (SE) is guided by a specification-based oracle that encodes how an implementation is expected to behave.

As discussed in \Cref{sec:paper2-limitations}, the effectiveness of specification-based oracles fundamentally depends on the completeness and precision of the requirements extracted from the protocol specification.
Violations can only be detected for behaviors that are explicitly captured by monitors, and any requirement that is omitted, misunderstood, or encoded imprecisely remains unchecked.
This dependence is particularly challenging for complex protocols such as DTLS and QUIC, whose specifications are lengthy, intricate, and in some cases ambiguous.
Even when SE explores a wide range of inputs, the analysis remains inherently incomplete, as it is constrained by the coverage of the encoded oracle.

In addition, not all protocol behaviors are governed by strict, explicitly stated rules.
Many aspects of protocol behavior arise from underspecified parts of the specification, where multiple implementation choices may be considered compliant.
Accurately capturing the full space of such permitted behaviors in the form of executable monitors is often difficult.
As a result, implementations may exhibit divergent behaviors that do not violate any explicitly stated requirement and therefore remain undetected by specification-based checking, even though such discrepancies may indicate subtle inconsistencies, corner cases, or potential interoperability issues.

This chapter considers protocol testing scenarios that arise precisely in the presence of such incomplete or underspecified oracles.
Rather than checking an implementation against an explicit set of specification-derived requirements, we adopt a comparative perspective and analyze how multiple implementations behave when exposed to the same inputs.
The underlying assumption is that independently developed implementations of the same protocol are expected to exhibit equivalent externally observable behavior when processing identical packet sequences.
Deviations from this common behavior may indicate implementation errors, inconsistent interpretations of the specification, or behaviors that are not explicitly characterized by the specification.

To explore such discrepancies systematically, this chapter introduces a differential SE technique for network protocol implementations.
The technique symbolically executes multiple implementations side by side on shared symbolic packet inputs and observes their externally visible behavior.
When SE produces execution paths whose responses differ across implementations for a common input, the subsequent analysis reports a behavioral discrepancy and extracts the concrete packet values that induce it.
At this stage, no judgment is made about which behavior is correct.
The resulting packet sequences are instead used as test cases whose effects are analyzed with reference to the protocol specification to determine whether the discrepancy corresponds to a violation, an underspecified behavior, or a permissible implementation choice.

The differential technique builds directly on the monitor-based infrastructure introduced in \Cref{chapter:paper2}.
The same test harness, packet parsing infrastructure, and SE setup are reused.
The key difference lies in how observations are interpreted.
Monitors are no longer used to enforce specification requirements, but instead to observe packet exchanges and record execution outcomes that are later analyzed to identify behavioral discrepancies across implementations.
This reuse allows differential SE to be integrated into the existing framework while addressing classes of protocol behaviors that are difficult to capture using specification-based oracles alone.
\section{Technique}
\subsection{Creating a Test Harness}
Differential SE is performed within a test harness that orchestrates a predefined interaction between a protocol implementation under test and another protocol party.
As in the monitor-based technique, the purpose of the test harness is to provide a controlled and deterministic execution environment.
In the differential setting, the harness additionally ensures that multiple implementations are exposed to equivalent inputs.

For each system under test (SUT), a separate test harness is constructed.
Each harness connects two components.
The first component is the SUT, which is the protocol implementation whose behavior is being analyzed.
The second component is another protocol party, which we refer to as the \emph{facilitator}.
The facilitator provides the inputs that drive the protocol interaction forward.
Depending on the protocol and the role of the SUT, this interaction may involve completing a handshake or, more generally, exchanging a predefined sequence of packets between protocol parties.

The facilitator implements the complementary protocol role to that of the SUT.
For example, in a client/server protocol such as DTLS, if the SUTs are server implementations, the facilitator is a client implementation.
For a fixed protocol role of the SUT, the same facilitator implementation is used across all tested SUTs.
Although a separate harness instance is constructed for each SUT, these harnesses share the same structure and differ only in the concrete SUT that they execute.
Reusing the same facilitator ensures that all SUTs are subjected to identical inputs and interaction patterns, and that any observed behavioral differences arise from differences in the SUT implementations rather than from variations in their peers.

To illustrate this setup, we use a running example based on the sequence number handling requirement discussed in \Cref{chapter:paper1}.
Suppose we wish to compare the behavior of the DTLS server implementations of Mbed TLS and wolfSSL.
Applying differential SE requires constructing two test harnesses, one for each server implementation.
Both harnesses use the same DTLS client implementation as the facilitator, for example a client from the Mbed TLS library, although other choices such as OpenSSL are equally possible.
Each facilitator drives the interaction by sending the same sequence of handshake messages to its corresponding server.
\Cref{fig:test-harness-arch} shows the architecture of one of the test harnesses used in this running example.
%
\begin{figure}[!ht]
	\centering
	\includegraphics[width=0.7\columnwidth]{figures/paper3:test\_harness.pdf}
	\caption{Architecture of the test harness used in the running example to test a DTLS server implementation.}
	\label{fig:test-harness-arch}
\end{figure}
%
\subsection{Running Symbolic Execution}
Once the test harness has been constructed, it can already be used to perform basic differential testing by executing multiple implementations on identical concrete inputs and comparing their observable behavior.
Such testing can reveal discrepancies for specific packet sequences, but it remains limited to concrete interactions that are manually selected.
Leveraging the power of SE extends this basic form of differential testing by systematically exploring variations in selected packet fields and by uncovering behaviors that arise only for specific input values and that may later cause discrepancies.

Differential SE is performed by executing each test harness under an SE engine.
As in the monitor-based technique, the harness executes in a controlled and deterministic environment.
However, rather than relying on assertions derived from specification requirements, SE is used to explore how different implementations respond to symbolic variations in protocol inputs.
Selected fields of packets generated by the facilitator are made symbolic, while all other fields retain their original concrete values.

The SE engine explores the code paths that are feasible for some assignment to the symbolic fields.
For each explored execution path, the engine records the externally observable outputs produced by the SUT.
These outputs include packets sent by the SUT or the absence of a response at the selected interaction point.
Technically, outputs are recorded as symbolic expressions that capture how they depend on the symbolic input fields.
This representation allows the analysis to reason about output behavior over a large input space without enumerating all concrete assignments explicitly.

Differential SE is conducted in a campaign-based manner.
For each pair of SUTs, multiple differential SE campaigns are executed.
Each campaign analyzes a specific protocol interaction and makes a specific set of fields symbolic.
For example, in DTLS, one campaign may focus on the handshake for a specific cipher suite and manipulate the record sequence number, while another campaign may target a different cipher suite.

Monitors play an important but fundamentally different role in the differential setting than in the specification-based technique introduced in \Cref{chapter:paper2}.
In differential SE, monitors do not encode protocol requirements and do not raise assertions.
Instead, they are used to observe packet exchanges and to record selected responses of each implementation during SE.
In particular, monitors record responses such as packets generated by the SUT or the absence of a response and associate these responses with the symbolic inputs that led to them.
This information is later used to compare the behavior of different implementations.
%
\begin{figure}[hb]
	\centering
		\begin{tikzpicture}[node distance=4.8cm]
			\node[state, initial, xshift=-.1cm] (c1) {\ST{Init}};
			\node[state, right of=c1] (c2) {\ST{Rcvd}};
			\node[state, right of=c2, xshift=-0.8cm] (c3) {\ST{End}};    
			
			\draw	
			(c1) edge[loop above, min distance=2em] node{$\frac{
					\edgelabel{P.source\, =\, Client\,}}{}$} (c1)
			(c1) edge node[above,  font=\scriptsize]{$
				\frac{
					\edgelabel{P.source\, =\, Client}}{
					\begin{aligned}
						&\edgelabel{make\_symbolic(\&P.sequence\_number);}
				\end{aligned}}$} (c2)
			(c2) edge[bend left=20] node[above]{$\frac{\edgelabel{P.source\, =\, Server}}{\edgelabel{record\_response(P)}}$} (c3)
			(c2) edge[bend right=20] node[below]{$\frac{\edgelabel{P.source\, =\, Client}}{}$} (c3);
		\end{tikzpicture}
	\caption{Monitor recording a DTLS server's responses for different values of \rfc{sequence\_number}.}
	\label{fig:monitor:dtls:sequence}
\end{figure}
%

In the running example, we apply this technique to compare the DTLS server implementations of Mbed TLS and wolfSSL.
To make the role of monitors in the differential setting concrete, \Cref{fig:monitor:dtls:sequence} shows the monitor used in the running example to observe DTLS server behavior for different values of the \rfc{sequence\_number} field.
The monitor is invoked on every packet exchanged during the interaction.
In the starting state \ST{Init}, it passively observes packets sent by the client.
At a nondeterministically chosen client packet, the monitor makes the \rfc{sequence\_number} field symbolic and transitions to the state \ST{Rcvd}.
From this point onward, the monitor waits for the next packet sent by the server.
When such a packet is observed, the monitor records the server's response using the \texttt{record\_response} action and transitions to the terminal state \ST{End}.
If another client packet is observed instead, the monitor assumes that the server has generated no response and transitions to \ST{End}.
Reaching the state \ST{End} terminates the current path and prompts the SE engine to record the corresponding path condition.

Running SE with such a monitor produces, for each SUT, a set of explored paths.
For every explored path, the SE engine records the corresponding path condition together with the selected interaction point and the response observed at that point.
To ensure that responses are comparable across SUTs, we associate each symbolic field with the position of the packet in the interaction, thereby fixing the interaction point at which responses are observed.
This position, which we refer to as the \emph{index}, identifies the interaction point at which the symbolic field was introduced.

The \emph{index} is not tracked explicitly by the monitor.
Instead, the underlying infrastructure assigns a concrete counter value to each packet processed during the interaction and introduces a symbolic variable that selects the packet on which the monitor will act.
This mechanism is the same as the one used in the monitor-based technique described in \Cref{chapter:paper2}.
Both the selected index and the symbolic field become part of the path condition recorded by SE.

Finally, to limit the number of explored paths and to reduce the manual effort required for analysis, each differential SE campaign employs a single monitor.
Such a monitor makes exactly one instance of a field symbolic, while all other instances of the same field in the interaction remain fixed to their concrete values.
This design isolates the effect of a single symbolic variation at a well-defined interaction point and keeps the analysis manageable.

\subsection{Finding and Analyzing Discrepancies}\label{sec:meth-discrepancies}
Given two SUTs, denoted $\mathtt{I}$ and $\mathtt{I'}$, the goal of differential SE is to identify common inputs that cause the two SUTs to produce different responses while being in the same interaction point.
Discrepancy detection is performed as a separate, offline analysis phase that operates on the results produced by SE.
The process consists of three main steps, which we describe below.

\paragraph{Deriving Path-Output Sets.}
For each SUT, SE explores a set of execution paths.
For every explored path, the SE engine records the path condition together with the response observed at the selected interaction point.
We represent each such record as a \emph{Path-Output} pair of the form $\langle pc, o \rangle$, where $pc$ denotes the path condition and $o$ denotes the recorded response.
Let $S_I$ denote the set of Path-Output pairs generated by the SE campaign for implementation $\mathtt{I}$.

\paragraph{Deriving Discrepancies Between Implementations.}
Having obtained the sets $S_I$ and $S_{I'}$ for two implementations $\mathtt{I}$ and $\mathtt{I'}$, we compare these sets to determine whether there exists a common input under which the two implementations produce different responses.
Because each Path-Output pair is associated with a specific packet index, comparisons are restricted to pairs that correspond to the same interaction point, ensuring that both SUTs are observed in the same interaction point.

If a concrete input can cause $\mathtt{I}$ and $\mathtt{I'}$ to produce different responses, then there must exist a Path-Output pair $\langle pc, o \rangle \in S_I$ and a Path-Output pair $\langle pc', o' \rangle \in S_{I'}$ such that both path conditions are satisfied by the same input and the resulting responses differ.
This condition can be expressed as the satisfiability of the following formula:
\begin{gather}
	\mathtt{(pc \, \land pc') \, \land \, (o \, \neq \, o')}\label{eq:discrepancy}
\end{gather}

The conjunct $\mathtt{(pc \land pc')}$ ensures that a common assignment to the symbolic input fields satisfies both path conditions, while the conjunct $\mathtt{(o \neq o')}$ ensures that this assignment leads to different responses.
For each pair of Path-Output pairs $\langle pc, o \rangle$ and $\langle pc', o' \rangle$ with the same packet index, we check the satisfiability of Formula \ref{eq:discrepancy} using an SMT solver.
If the formula is satisfiable, we ask the solver to produce a valuation for the symbolic input field and the corresponding packet index.
Each such valuation, together with the differing responses $o$ and $o'$, constitutes a \emph{discrepancy witness}.
Our technique produces a witness for every pair of Path-Output pairs that satisfies Formula \ref{eq:discrepancy}.

\paragraph{Analyzing Discrepancies for Bugs.}
For each discrepancy witness, the protocol specification is consulted to determine whether the observed behaviors are compliant.
At this stage, the RFC serves as a reference for interpretation rather than as an executable oracle.
If one of the implementations is determined to be non-compliant, we validate the discrepancy by executing the test harness on the unmodified implementation.
The witness provides both the value of the symbolic field and the index of the packet in which it appears.
During validation, the corresponding packet is modified to contain the witness value, and the interaction is replayed to confirm that the divergent behavior can be reproduced.
If confirmed, the non-conformance is reported to the developers.

We conclude this section by instantiating the above procedure for the running example involving the \rfc{sequence\_number} field.
Symbolic execution produces 37 Path-Output pairs for Mbed TLS and 45 Path-Output pairs for wolfSSL.
Although these numbers may appear large, all steps up to discrepancy interpretation can be automated.
Instantiating Formula \ref{eq:discrepancy} for these Path-Output pairs yields three satisfying assignments, each forming a discrepancy witness.
Each witness contains a concrete value for \rfc{sequence\_number}, the index of the packet containing it, and the corresponding responses produced by the two SUTs.

One of the witnesses indicates that, in the initial interaction point, upon receiving a first \emph{ClientHello} packet with a \rfc{sequence\_number} slightly larger than $2^{44}$, Mbed TLS responds with \emph{HelloVerifyRequest}, indicating acceptance of the packet, while wolfSSL produces no response and appears to silently discard it.
Consulting the DTLS specification~\cite{RFC6347} reveals that Mbed TLS handles this packet correctly, whereas wolfSSL incorrectly rejects a packet whose sequence number is large but still valid.
This discrepancy exposes a bug in wolfSSL that can lead to interoperability problems.

\subsection{Selecting Fields for Symbolic Execution}
The choice of which protocol fields to make symbolic, and which fields of response packets to record, has a significant impact on both the effectiveness and the practicality of differential SE.
Making too many fields symbolic, or recording overly detailed responses, can quickly lead to an explosion in the number of explored paths and identified discrepancies, many of which may not correspond to meaningful specification violations.
In this section, we outline general guidelines for selecting input and output fields, following the technique adopted in our implementation.

\paragraph{Selecting Fields to Make Symbolic in Input Packets}
In principle, our technique can be applied to any protocol field, including fields such as \rfc{sequence\_number}, whose incorrect handling may have serious security or functional implications.
In practice, however, performance considerations and limitations of SE necessitate a more selective approach.
In particular, fields that are fed into cryptographic operations are often difficult to handle symbolically.
Symbolic execution engines typically struggle with cryptographic functions, and accessing encrypted fields within monitors may require additional infrastructure, such as packet decryption.
To avoid these complications, one option is to restrict symbolic manipulation to fields that are transmitted in plaintext.
In DTLS, for example, the \rfc{sequence\_number} field is not encrypted and can therefore be accessed and manipulated without decrypting records.

Another option is to disable encryption altogether.
This can be achieved either through configuration settings supported by the implementation under test or by applying small, targeted modifications to the SUT.
In our DTLS experiments, we disabled encryption using configuration options provided by the tested implementations.
\paragraph{Selecting Fields to Record in Response Packets}\label{sec:outputfieldselection}
Similarly to input fields, one could in principle record entire response packets in order to maximize the number of discrepancies that can be identified.
In practice, however, recording all response fields tends to produce a large number of discrepancies that are unlikely to correspond to bugs and that significantly increase the manual effort required for analysis.
For this reason, our technique selectively excludes certain classes of response fields from consideration.

The first class consists of fields that contain nonces or other randomly generated values.
Recording such fields would almost always lead to trivial discrepancies, since independently generated random values are expected to differ across implementations.
For example, we do not record the \rfc{random} field in the DTLS \emph{ServerHello} message, as this field contains a nonce.

The second class consists of fields whose values primarily reflect differences in supported features or configuration options across implementations.
Such differences are often unavoidable and do not necessarily indicate incorrect behavior.
Recording these fields would therefore result in discrepancies that are unlikely to be representative of actual nonconformance issues.
As an example, we do not record the \rfc{cipher\_suites} field in the DTLS \emph{ServerHello} message, since its value depends on the set of cryptographic algorithms supported by the implementation, which varies widely across SUTs.

The third class consists of fields whose values depend on fields belonging to the first or second class.
The same reasons for excluding those fields apply here as well.
For instance, the DTLS \emph{Finished} message contains the \rfc{verify\_data} field, which is computed as a hash over several protocol fields, including \rfc{random} and \rfc{cipher\_suites}.
We therefore exclude this field from analysis.

In our implementation, the selection and recording of response fields is handled by a dedicated output-capture component.
This component is responsible for extracting the chosen fields from response packets and associating them with the corresponding path conditions.
We describe this component in the next section.
\section{Implementation}
This section describes how the differential SE technique introduced in the previous sections is realized in practice.
The implementation builds directly on the monitor-based infrastructure presented in \Cref{chapter:paper2}.
In particular, the same packet interception mechanism, monitor invocation model, and protocol-specific parsing infrastructure are reused.
The main differences lie in how monitors are used, how protocol responses are recorded, and how the results of SE are analyzed to identify discrepancies across implementations.

We first describe the construction of the test harness.
We then outline the extensions to the SE setup that enable the generation of Path-Output pairs.
Finally, we describe how discrepancies are identified and discuss optional adaptations to the SUT that improve SE performance.

\subsection{Creating a Test Harness}
For each SUT, SE is performed within a dedicated test harness.
As in the monitor-based technique, the harness orchestrates a predefined interaction between the SUT and another protocol party.
The role of the harness is to provide a controlled and deterministic execution environment in which packet exchanges are fully mediated.

Each test harness consists of two components.
The first component is the SUT, which is the protocol implementation whose behavior is analyzed.
The second component is a peer implementation, referred to as the \emph{facilitator}, which drives the interaction by sending protocol messages to the SUT.
In client/server protocols such as DTLS, when the SUT is a server implementation, the facilitator acts as a client.

All components execute within a single Unix process.
This design choice is dictated by the SE engine used in our implementation, KLEE~\cite{KLEE@OSDI-08}, which supports only single-threaded execution.
Communication between the facilitator and the SUT is therefore realized using an in-memory packet exchange mechanism rather than operating system sockets.
As in the monitor-based technique, this is achieved by overriding standard POSIX socket APIs and redirecting packet transmission through shared queues.

Protocol implementations are typically distributed as libraries accompanied by example programs.
We construct test harnesses by adapting these example programs.
In the DTLS setting, we start from the sample client and server programs provided by each library and refactor them into facilitator and SUT components.
To ensure that all SUTs are exposed to identical inputs, the same facilitator implementation is used across all tested SUTs.
Only the SUT component varies between harnesses.
Although a separate harness binary is required for each library, a harness can typically be reused across multiple versions of the same implementation with little or no modification.

For DTLS, each harness executes a handshake between the facilitator and the SUT and then exchanges a small piece of application data.
Execution proceeds in a loop that terminates when the handshake has completed on both sides and exactly one record containing application data has been exchanged in each direction. 
The loop also terminates if a fatal alert is generated or if an iteration produces no new data.
To prevent unbounded execution, we additionally impose a fixed upper bound on the number of iterations.

\subsection{Augmentations to Symbolic Execution}
Our technique requires extending the SE setup so that protocol responses can be recorded and compared across implementations.
These extensions reuse large parts of the monitor-based infrastructure described in \Cref{chapter:paper2}, while adapting it to support this technique.
\subsubsection{Common Infrastructure}
As in the monitor-based technique, packet exchange between the facilitator and the SUT is mediated by a proxy that intercepts POSIX socket calls.
When executed under KLEE, the harness invokes custom implementations of socket functions such as \texttt{sendto} and \texttt{recvfrom}, which enqueue and dequeue packet payloads from in-memory queues.
Each direction of communication maintains a separate queue.

Whenever a packet is dequeued, it is first parsed into a structured representation using protocol-specific infrastructure.
A monitor is then invoked on this representation.
After the monitor completes its execution, the possibly modified packet is serialized back into its wire format and delivered to the receiving party.
This execution model is unchanged from the monitor-based framework and is reused without modification.

\subsubsection{Protocol-Specific Infrastructure}
In addition to the common infrastructure, our technique relies on several protocol-specific components.
These include a set of monitors, an output-capture function, and optional parser and serializer components.

\paragraph{Set of Monitors}
Monitors in the differential setting are considerably simpler than those used for specification-based testing.
Their role is not to evaluate protocol requirements or check assertions, but to select fields for symbolic exploration and to trigger the recording of responses.

Each monitor is implemented as a C function that operates on a structured packet representation.
When invoked, the monitor may make a specific packet field symbolic and then wait for a corresponding response from the SUT.
Only one instance of a field is made symbolic per SE campaign.
All other instances of the same field remain concrete.

\Cref{fig:dse-implementation:monitor} shows the C implementation of the monitor used in the running example to explore variations of the DTLS \rfc{sequence\_number} field.
The monitor maintains a simple control state.
In the initial state, it observes client packets and nondeterministically selects one packet on which to act.
When that packet is encountered, the \rfc{sequence\_number} field is made symbolic.
The monitor then waits for the next packet generated by the server and invokes the output-capture function (\texttt{capture\_dtls\_output}) to record the response.

As in the monitor-based framework, nondeterministic selection of the packet index is implemented by the infrastructure rather than explicitly encoded in the monitor.
Each packet processed during execution is assigned a concrete index, and a symbolic variable ranges over all possible packet indices.
The monitor transitions from its initial state only when the current packet index matches this symbolic value.
Both the selected index and the symbolic field value are included in the path condition generated by SE.

\begin{figure}[t]
	\begin{lstlisting}[language=C, numbers=left, xleftmargin=1em, frame=single,
		numberstyle=\tiny\color{blue}, showstringspaces=false,
		escapeinside={//(*@}{@*)}]
enum STATE { INIT, RCVD, END }; 

static STATE curr_state = INIT;

void sequence_number_diff_testing_server(RECORD *P) {
	if (curr_state == INIT && P.source == CLIENT) {
		kleener_make_symbolic(&P->sequence_number,
			sizeof(P->sequence_number), "sequence_number");
		curr_state = RCVD;
	}
	else if (curr_state == RCVD && P.source == SERVER) {
		capture_dtls_output(P);
		curr_state = END;
	}
	else if (curr_state == RCVD && P.source == CLIENT)
		curr_state = END;
	else if (curr_state == END) return;
}
	\end{lstlisting}
	\caption{The C implementation of the monitor of~\cref{fig:monitor:dtls:sequence}.}
	\label{fig:dse-implementation:monitor}
\end{figure}

\paragraph{Output-Capture Function}
Output recording is implemented via a dedicated output-capture function that is invoked by the monitor when a response from the SUT is observed.
For DTLS, this function inspects the parsed packet produced by the SUT, selects the response fields according to the guidelines in \Cref{sec:outputfieldselection}, and records them as part of the observed response.

When a recorded field depends on symbolic input, the output-capture function records the corresponding symbolic expression.
As a result, each explored execution path yields a Path-Output pair consisting of a path condition and a symbolic representation of the observed protocol response.
These Path-Output pairs form the basis for discrepancy detection in the subsequent analysis phase.

\paragraph{Parser and Serializer}
As in the monitor-based technique, optional parser and serializer components are used to translate between raw packet bytes and structured protocol representations.
These components simplify the implementation of monitors and output-capture functions by allowing direct access to protocol fields.
They are reused unchanged from the earlier framework.

\subsection{Finding and Analyzing Discrepancies}
Executing the test harness under SE produces a set of Path-Output pairs for each tested SUT.
Each pair consists of an SMT formula encoding the path condition and a symbolic representation of the response recorded by the output-capture function.

For each SE campaign and each tested SUT, KLEE produces one constraint file per explored path, containing the path condition, and a corresponding file that encodes the recorded response.
Both artifacts are expressed in SMT-LIBv2 format.
To identify discrepancies, we compare the Path-Output pairs generated for different SUTs using an offline analysis script.

For every pair of SUTs, the script iterates over pairs of paths that correspond to the same packet index and checks the satisfiability of the discrepancy condition represented by Formula~\ref{eq:discrepancy} using the SMT solver Z3.
If the condition is satisfiable, the solver produces a concrete assignment to the symbolic variables that causes the two SUTs to exhibit different responses.
This assignment, together with the differing responses, constitutes a discrepancy witness.

Because discrepancy witnesses are analyzed manually, the script applies a simple deduplication strategy.
A witness is reported only if no previously reported witness exists with the same packet index and the same pair of observed responses.
This filtering significantly reduces the number of discrepancies that must be inspected while preserving representative cases.
In our running example, this strategy reduces the number of reported discrepancies from ten to three.

\subsection{Optional SUT Adaptations to Speed Up SE}
Symbolic execution is substantially slower than native execution, and protocol implementations often include behaviors that further increase analysis complexity.
As in the monitor-based technique described in \Cref{chapter:paper2}, we apply a small set of optional adaptations to the SUT to improve performance without affecting the logical behavior under analysis.

In particular, we suppress time-based behavior such as retransmissions and timeouts, either through configuration options provided by the implementation or through minor code modifications when necessary.
We also avoid expensive cryptographic operations during SE.
For DTLS, this is achieved by configuring the SUT and facilitator to disable encryption after the handshake.


\subsection{Evaluation}
We evaluated our technique to answer four main questions.
First, how does the technique compare to the monitor-based, specification-driven approach in terms of bug-finding effectiveness.
Second, how does the required testing effort compare, and how can this effort be optimized.
Third, can the technique uncover previously unknown bugs in recent versions of DTLS implementations.
Finally, can the technique be used to detect regressions across implementation versions.

\paragraph{Comparison with Monitor-Based Specification Testing}
The monitor-based technique presented in \Cref{chapter:paper2} relies on explicitly encoded specification requirements to detect violations.
As a result, its effectiveness depends directly on the completeness and precision of the specification-based oracle.
Differential SE removes this dependency by comparing the externally observable behavior of multiple implementations under identical symbolic inputs.

In our evaluation, differential SE uncovered discrepancies that were not detected by the monitor-based approach.
Several of these discrepancies arose in parts of the protocol that are underspecified or loosely defined in the RFC, making them difficult to encode accurately as specification-based monitors.
While the monitor-based technique excels at detecting clear violations of explicitly stated requirements, differential SE proved more effective at revealing subtle semantic differences and corner cases that manifest only through comparison.
The two techniques are therefore complementary rather than redundant.

\paragraph{Testing Effort and Optimization}
Differential SE requires less manual effort to define protocol oracles, as it does not require extracting and encoding individual specification requirements.
Instead, the main manual effort lies in selecting protocol interactions, choosing fields to make symbolic, and interpreting the discrepancies produced by the analysis.

The symbolic exploration cost is comparable to that of the monitor-based technique, since both approaches rely on similar test harnesses and SE infrastructure.
However, differential analysis may initially produce a larger number of candidate discrepancies.
We mitigate this by restricting each SE campaign to a single symbolic field instance and by applying deduplication during discrepancy analysis.
These optimizations significantly reduce the number of discrepancies that require manual inspection and keep the overall testing effort manageable.

\paragraph{Detection of Previously Unknown Bugs}
We applied differential SE to several widely used DTLS server implementations.
The evaluation uncovered nine distinct bugs, five of which were previously unknown, including deviations from the protocol specification that can lead to interoperability issues.
Importantly, these bugs were found in up-to-date versions of the implementations, demonstrating that the technique remains effective even as implementations evolve.

Several of the discovered issues arise from subtle semantic differences in how implementations interpret and enforce protocol rules.
Such issues are difficult to capture using specification-based testing alone, particularly when the relevant aspects of the specification are underspecified or permit multiple interpretations.
By systematically comparing behaviors across implementations, differential SE exposes these discrepancies and provides concrete witnesses for further analysis.

\paragraph{Regression Detection}
Differential SE can also be used to detect regressions by comparing different versions of the same implementation.
By treating two versions of a library as separate SUTs and executing identical SE campaigns, the technique can reveal behavioral changes introduced by code modifications.
This allows developers to identify unintended regressions or semantic changes, even when such changes do not violate explicitly encoded specification requirements.

Overall, the evaluation demonstrates that differential SE is an effective and practical complement to specification-based SE.
By removing reliance on specification-based oracles, the technique enables the discovery of protocol nonconformances and semantic discrepancies that are difficult to detect using explicit requirement checks alone.
A detailed description of the experimental setup, tested implementations and versions, and the full set of discovered issues is provided in Paper \URoman{3}.

Building on the techniques presented in earlier chapters, this chapter explored differential SE as a means of testing network protocol implementations without relying on specification-based oracles.
By comparing the externally observable behavior of multiple implementations under symbolic inputs, the technique complements the requirement-driven and monitor-based approaches presented earlier in this thesis.

Symbolic execution enables systematic exploration of protocol behavior and precise reasoning about corner cases.
Fuzzing offers a complementary perspective by enabling lightweight, scalable exploration of large input spaces.
The final chapter brings these approaches together by combining fuzzing and symbolic execution to test protocol requirements in the context of IoT protocols.
Focusing on the Constrained Application Protocol (CoAP), it demonstrates how these techniques can be applied effectively in resource-constrained environments.

    \chapter{Two Paths to Conformance Testing: Fuzzing and Symbolic Execution}\label{chapter:paper4}
Network protocol conformance can be assessed through different testing techniques, each offering distinct strengths and limitations.
The previous chapters have focused on symbolic execution (SE) as a principled technique for systematically exploring protocol behavior and for reasoning about requirement violations that arise only under narrowly constrained conditions.
By progressively extending SE with requirement-driven oracles, external monitors, and differential analysis, we demonstrated how it can be applied effectively to complex, stateful protocols such as DTLS and QUIC.

However, SE is not the only technique used in practice to test protocol implementations.
Coverage-guided fuzzing has become a widely adopted approach for testing networked systems, particularly in settings where scalability and low analysis overhead are important.
Fuzzers excel at exercising large input spaces quickly and can expose robustness issues as well as logical flaws that manifest through externally observable behavior.
At the same time, fuzzing is not typically employed to check protocol requirements or to systematically target specific semantic conditions.

This chapter applies the requirement-centric perspective developed in earlier chapters to both techniques and compares their effectiveness when applied to the same requirements and implementations.
In particular, it examines how fuzzing and SE differ in their ability to expose requirement violations in the context of an IoT protocol.

\section{Technique}
\subsection{Test Harness Creation}
Testing protocol conformance requires a controlled execution environment in which interactions with a protocol entity can be reproduced and manipulated systematically.
As in the previous chapters, this is achieved through a dedicated test harness that mediates communication between the system under test (SUT) and its peer and that provides full control over the messages exchanged during a protocol interaction.

The construction of the harness begins by capturing a representative interaction with the protocol entity.
For CoAP, this interaction consists of a sequence of request and response messages exchanged between a client and a server.
Packet capture is performed by configuring sample programs provided with the SUT to execute a valid interaction and recording the packets sent to the protocol entity using standard network analysis tools such as TCPdump or Wireshark.
The resulting packet sequence serves as a concrete reference interaction that can be replayed during testing.

During testing, the harness advances the interaction step by step by replaying the pre-captured packets.
At each step, the harness invokes the appropriate API of the SUT to deliver an incoming message and to obtain the corresponding response.
Responses produced by the SUT are recorded and used by the harness to determine how the interaction proceeds.
This execution model ensures that the protocol logic exercised during testing follows the same high-level interaction pattern as the captured execution, while allowing individual message fields to be varied or constrained as required by the testing technique.

The test harness should be designed to support all protocol requirements considered for testing.
Depending on the nature of the requirements, this may involve executing a single interaction that exercises multiple protocol features or running several distinct interactions that target different parts of the protocol.
In either case, the harness provides a uniform mechanism for delivering inputs, observing outputs, and controlling progress through the protocol interaction.

A key design goal of the harness is to enable both fuzzing and SE with minimal adaptation.
The same harness structure is used in both settings, with differences limited to how input messages are generated and how responses are analyzed.
For fuzzing, the harness supplies mutated variants of the captured packets, while for SE it designates selected fields as symbolic and relies on the symbolic executor to explore feasible variations.
By sharing a common harness, both techniques are applied to identical interaction patterns and protocol requirements, enabling a meaningful comparison of their effectiveness in exposing requirement violations.

\subsection{Encoding Protocol Requirements as Assertions}
Network protocol specifications define requirements not only over individual messages but also over sequences of messages exchanged between protocol entities.
For CoAP, such requirements are specified across several RFC documents, most notably RFC~7252~\cite{RFC7252} and RFC~7959~\cite{RFC7959}.
Once relevant requirements have been extracted from these specifications, they can be translated into executable assertions that are evaluated during a protocol interaction.

In this work, protocol requirements are encoded as assertions over the packets exchanged during an interaction.
These assertions are inserted at points in the SUT where response messages are generated, allowing the observed behavior of the implementation to be checked against the specification at runtime.
In the CoAP implementations considered, assertions were inserted into the functions responsible for transmitting response messages.
These functions already operate on structured representations of outgoing packets, which allows assertions to directly inspect response fields.
Fields of the corresponding incoming packet are accessed through a global pointer that is maintained by the test harness and updated at each interaction step.

We illustrate this approach by presenting the encoding of two representative CoAP requirements.


\paragraph{Message ID Echo.}
The CoAP specification mandates that the Message ID of certain request messages must be echoed in the corresponding response.
RFC~7252~\cite[p.~24]{RFC7252} states:
%
\rfcdisplayquote{%
	The Message ID is a 16-bit unsigned integer that is generated by the sender of a Confirmable or Non-confirmable message and included in the CoAP header.
	The Message ID MUST be echoed in the Acknowledgement or Reset message by the recipient.}%
%
Let $\rfcm{p_{in}}$ denote the parsed incoming packet and $\rfcm{p_{out}}$ the response packet generated by the SUT.
The requirement can be encoded by the following assertion, which is evaluated when the response packet is constructed:
%
\begin{equation*}
	\begin{split}
		& \scriptstyle \mathtt{assert} \, (((\rfcm{p_{in}.type}\, =\, \rfcm{CON} \, \lor \, \rfcm{p_{in}.type}\, =\, \rfcm{NON}) \, \land \, \\[-5pt]
		& \scriptstyle \quad\quad \,\,\,\,\,\,\,(\rfcm{p_{out}.type}\, =\, \rfcm{ACK} \, \lor \, \rfcm{p_{out}.type}\, =\, \rfcm{RST})) \implies \\[-5pt]
		& \scriptstyle \qquad\qquad\qquad \rfcm{p_{in}.message\_id}\, =\, \rfcm{p_{out}.message\_id})
	\end{split}
\end{equation*}
%
This assertion ensures that whenever the SUT generates an acknowledgement or reset in response to a Confirmable or Non-confirmable request, the Message ID is correctly echoed.
During testing, the assertion is evaluated at each interaction step, and any violation immediately indicates nonconformance with the specification.

\paragraph{Further Request Block Size.}
Some CoAP requirements relate the handling of a request to information conveyed in earlier responses.
RFC~7959~\cite[p.~14]{RFC7959} specifies such a relationship for block-wise transfers:
%
\rfcdisplayquote{%
	In response to a request with a payload (e.g., a PUT or POST transfer), the block size given in the Block1 Option indicates the
	block size preference of the server for this resource.
	\dots
	the client SHOULD heed the preference indicated and, for all further
	blocks, use the block size preferred by the server or a smaller one.}%
%

To encode this requirement, the assertion logic explicitly maintains state across packets.
A boolean variable $\rfcm{first\_response}$ indicates whether the current response is the first generated by the server, and a variable $\rfcm{preferred\_size}$ stores the block size advertised by the server.
The requirement is checked using the following logic:
%
\begin{equation*}
	\begin{split}
		& \scriptstyle \mathtt{if} \, (\, \rfcm{first\_response} \, \land \, \rfcm{p_{in}.code} \, \in \, \rfcm{REQ\_CODES})\\[-5pt]
		& \scriptstyle \quad \rfcm{preferred\_size} \, := \, \rfcm{BlockSize (p_{out})}\\[-5pt]
		& \scriptstyle \quad \rfcm{first\_response} \, := \, \mathtt{false} \\[-5pt]
		& \scriptstyle \mathtt{else \, if} (\neg \, \rfcm{first\_response})\\[-5pt]
		& \scriptstyle \quad \mathtt{assert} ((\rfcm{BlockSize (p_{in})} \, > \, \rfcm{preferred\_size}) \, \implies \\[-5pt]
		& \scriptstyle \qquad\qquad\qquad \rfcm{p_{out}.code} \, \in \,  \rfcm{ERR\_CODES})
	\end{split}
\end{equation*}
%
The assertion is evaluated only after the server’s preferred block size has been established by the first relevant response.

Here, \rfc{REQ\_CODES} represents the set of CoAP request codes.
The function $\rfcm{BlockSize}$ extracts the block size specified by the Block1 option of a packet, and \rfc{ERR\_CODES} denotes the set of CoAP error response codes.
The assertion checks that if a client continues to send request blocks larger than the server's advertised preference, the server responds with an error.

While the RFC does not prescribe a specific error response for this case, requiring an error response provides a reasonable approximation of the intended behavior and allows violations to be detected systematically.

Once the SUT has been extended with such assertions, they serve as executable representations of protocol requirements.
Both fuzzing and SE can then be applied to explore whether the implementation can be driven into executions that violate these assertions, thereby revealing instances of protocol nonconformance.

\subsection{Testing Requirements Using Fuzzing}
While coverage-guided fuzzing is traditionally used to discover crashes and robustness issues, it can be adapted to support requirement-driven testing when combined with explicit assertions.
In this work, we use fuzzing not as a robustness testing technique, but as a mechanism for exploring executions that may violate protocol requirements encoded as runtime assertions.
Standard fuzzing infrastructure can be leveraged for this purpose without modifying the fuzzer itself, by generating inputs that may trigger violations of the encoded assertions.
In our evaluation, we use AFL++~\cite{AFLplusplus-Woot20}, a state-of-the-art fuzzer that incorporates techniques from recent research and provides efficient feedback-driven input generation.

Fuzzing is performed by executing the test harness described earlier.
As initial seeds, we use valid, pre-captured packets obtained from protocol interactions executed using the sample programs provided with each implementation.
These packets serve as well-formed starting points from which the fuzzer generates mutated inputs.
During a fuzzing campaign, the fuzzer repeatedly executes the harness and supplies mutated variations of the incoming packets to the SUT.
All encoded assertions are evaluated during each execution, and any violation is immediately detected.

Some protocol requirements, particularly those related to block-wise transfer in CoAP, depend on the processing of multiple packets and therefore require that protocol state be built across several interaction steps before an assertion can be checked.
Standard fuzzers are not designed to reason about such stateful interactions.
To accommodate these requirements, we represent the packets of a complete interaction as a single contiguous input buffer, with each packet placed at a fixed offset.
During execution, the harness delivers packets to the SUT by reading from the appropriate offsets within this buffer.
Since the fuzzer is free to arbitrarily mutate the input buffer, packet boundaries are not guaranteed to remain intact.
As a result, some mutated inputs may cause packet boundaries to shift or packets to become malformed.
In such cases, the harness may fail to parse the input correctly or the SUT may immediately reject the input, and these executions do not contribute to requirement checking.
When packet parsing succeeds, however, this design allows the fuzzer to mutate a single input buffer while still exercising a multi-packet protocol interaction.
As a result, state-dependent requirements can be tested using standard fuzzing infrastructure without modifying the fuzzer itself.

Taking into account the inherent randomness of fuzzing, campaigns are repeated multiple times to increase confidence in their ability to expose violations.
When the fuzzer triggers an assertion violation, it stores the corresponding mutated input together with metadata such as the time of discovery and the number of mutations applied.
These concrete test cases can then be analyzed to determine which packet fields and values caused the requirement violation and to reproduce the behavior on the unmodified implementation.

Overall, fuzzing provides a lightweight and scalable means of exploring large portions of the input space and of exercising multiple requirements simultaneously.
However, its effectiveness depends on the fuzzer's ability to reach inputs that satisfy the specific conditions under which assertions fail.
In the next subsection, we contrast this behavior with SE, which reasons explicitly about such conditions and can systematically target individual requirements.

\subsection{Testing Requirements Using Symbolic Execution}
The approach for testing requirements using SE employed in this work closely follows the requirement-driven SE technique introduced in Chapter~\ref{chapter:paper1}, using the same test harness structure and execution model.
The main differences in this work concern the protocol under test and the specific requirements considered, rather than the underlying SE methodology.
To test protocol requirements using SE, we employ KLEE~\cite{KLEE@OSDI-08}, which we use as the SE engine throughout this thesis.
As in the fuzzing-based approach, SE is performed by executing the test harness described earlier.
However, unlike fuzzing, SE targets individual requirements explicitly and is therefore performed in separate runs, with each run checking exactly one requirement.

In each SE run, KLEE executes the test harness using a pre-captured sequence of valid packets that define the interaction under test.
Before a packet is processed by the protocol party, the fields relevant to the requirement being checked are designated as symbolic, while all other fields remain concrete.
This selective symbolic treatment restricts exploration to the part of the input space that is relevant for the requirement and avoids unnecessary path explosion.

As in the earlier chapters, packets are parsed into structured representations to allow symbolic manipulation of individual protocol fields, and are serialized back into byte buffers before being processed by the protocol implementation.

As an example, to test the Message ID Echo requirement, the \rfc{message\_id} and \rfc{type} fields of the incoming packet are made symbolic.
KLEE then explores all feasible execution paths that arise from different assignments to these fields.
If there exists an execution path along which the corresponding assertion is violated, KLEE reports the violation and produces a concrete assignment for the symbolic fields that triggers it.
This assignment constitutes a concrete test case that demonstrates nonconformance with the requirement.

In this setting, a test case corresponds to the specific values assigned to the symbolic protocol fields, rather than to an entire packet sequence.
These values can be substituted back into the original packet structure and replayed against the unmodified implementation to reproduce the violation.
Symbolic execution therefore provides a systematic means of determining whether a requirement can be violated at all and, if so, of generating precise witnesses that expose the violation.

\section{Evaluation}
The purpose of the evaluation is to assess how fuzzing and SE behave in practice when applied to the same protocol requirements within a shared testing framework.
The evaluation focused on a set of representative protocol requirements extracted from the CoAP specification and its extensions.

Overall, our findings demonstrate that both fuzzing and SE are effective at exposing protocol nonconformances when requirements are encoded explicitly as assertions.
Across the evaluated requirements, both fuzzing and SE identified the same set of nine distinct nonconformances in the tested implementations.
All identified issues were reported to the respective developers, and most have been fixed at the time of writing.
These results indicate that assertion-based testing is capable of uncovering concrete and actionable specification violations in real-world CoAP implementations.

In terms of effectiveness at detecting violations, we did not observe a substantial difference between the two techniques.
For the evaluated requirements, fuzzing and SE were able to identify violations for the same subsets of requirements in each implementation.
Moreover, violations were typically discovered quickly.
Fuzzing consistently generated inputs that triggered assertion failures within seconds for most requirements, while SE also produced violating inputs within comparable time frames in the majority of cases.
This suggests that, when requirement violations are reachable through relatively shallow protocol interactions, both techniques are well suited to exposing them.

The main distinction between fuzzing and SE becomes apparent when considering requirements that an implementation respects.
Fuzzing is inherently probabilistic and cannot provide evidence that a requirement is upheld beyond the absence of observed violations within a given time budget.
Symbolic execution, by contrast, explores program paths systematically and can, in principle, establish that no execution violating a requirement exists within the explored search space.
In practice, this ability is limited by path explosion and solver performance.
For some requirements, SE was able to complete exploration and thus provide stronger assurance of correctness.
For others, exploration did not complete within the imposed time limits.

Taken together, these observations suggest a complementary relationship between the two techniques.
Fuzzing offers a lightweight and scalable way to uncover specification violations quickly, making it well suited for early testing.
Symbolic execution, while more computationally expensive, provides deeper insight into the conditions under which violations occur and can offer stronger guarantees when exploration completes.
Our experience indicates that combining both techniques within a single assertion-based framework allows developers to benefit from the strengths of each while mitigating their respective limitations.
A detailed description of the experimental setup, the evaluated implementations, and the complete set of discovered issues is provided in Paper \URoman{4}.


    \chapter{Related Work}\label{chapter:related}
This chapter reviews related work on techniques for testing protocol implementations, focusing on methods that test network protocol implementations.

\section{Specification-Based Protocol Testing}
Ensuring that network protocol implementations conform to their specifications has long been an important goal, motivated by the fact that even subtle deviations can lead to interoperability issues or security vulnerabilities.
Traditional approaches for testing conformance are typically specification-driven and rely on abstract models or explicitly encoded properties derived from specification documents such as RFCs.

Model-based Testing (MBT) is a well-established approach in this space, where an abstract behavioral model of the protocol is either manually constructed or inferred through learning techniques~\cite{MBT-reactive@2005,Utting:mbt@2006,Tappler@ICST-2017}.
This model is then used to systematically generate test inputs, which are applied to an implementation, usually in a black-box setting.
A wide range of modeling formalisms and tooling support has been proposed, including automata-based specifications and domain-specific languages~\cite{Veanes:specexplorer,McMillan:sigcomm19}.

Recent work has demonstrated the applicability of MBT to realistic networked and distributed systems.
Li \etal~\cite{LiMBT@ISSTA21} derive executable tests for networked applications from formal models, while MBT has also been applied to IoT protocols such as BLE~\cite{MBT-IoT@FMICS22} and MQTT~\cite{MBT-MQTT@KBSE20}, uncovering conformance issues in implementations of these protocols.

While MBT can provide strong guarantees when accurate models are available, constructing and maintaining such models for complex evolving protocols remains labor intensive.

Property-based testing (PBT) offers a related but often lighter-weight alternative~\cite{ArtsQuviq@Erlang-06,PBT-SensorNetworks@SECON-15,TPBT@ISSTA-17, TPBT@ICST-18}.
Instead of a full behavioral model, PBT relies on properties that inputs and outputs should satisfy, combined with automated input generation.
The high-level abstraction and built-in generators provided by PBT frameworks often make them easier to apply than MBT.
However, as with MBT, the effectiveness of PBT depends on the completeness and precision of the manually specified properties.

Both MBT and PBT are commonly applied in black-box testing scenarios.
As a result, they are limited in their ability to guide exploration toward specific internal program states or code paths that may expose subtle requirement violations.
In contrast, white-box techniques can leverage internal program structure to search more directly for inputs that trigger specification violations, especially in the presence of complex control flow and stateful behavior.

\section{Fuzzing for Network Protocol Implementations}
Fuzz testing has established itself as one of the most effective techniques for uncovering software defects by repeatedly executing a program on mutated inputs and observing unintended behavior, particularly crashes and memory errors.
Modern coverage-guided greybox fuzzers, such as AFL~\cite{AFL} and its successor AFL++~\cite{AFLplusplus-Woot20}, use lightweight instrumentation to guide input generation toward previously unexplored branches and have achieved considerable success across a wide range of software systems.

Fuzzing has also been applied to network and IoT protocol implementations.
For example, fuzzers have been used to test IPv6 network stacks in constrained environments~\cite{SoManyFuzzers@ASE-22}, as well as protocol implementations such as CoAP~\cite{liljedahl19,melo-17,zeng-20}.

A key challenge for fuzzing protocol implementations is statefulness.
Many protocols require specific sequences of messages to reach deep states in the implementation, and naive fuzzing often fails to explore such interactions.
To address this challenge, several extensions to greybox fuzzing have been proposed, including AFLNET~\cite{AFLNET@ICST-20}, StateAFL~\cite{StateAFL@EmpirSoftEng-22}, SGFuzz~\cite{SGFuzz@USENIX-22}, SnapFuzz~\cite{SnapFuzz@ISSTA-22} and LibAFLstar~\cite{LibAFLStar@2025}.
These approaches incorporate protocol state awareness, for instance by tracking message sequences or using lightweight protocol models, in order to improve coverage in stateful systems.
Empirical evaluations show that such techniques significantly outperform traditional greybox fuzzers on network protocol implementations, enabling deeper exploration of protocol logic and leading to the discovery of previously unknown bugs across a range of stateful protocols.
Nevertheless, their effectiveness is primarily measured in terms of coverage and crash discovery, and they do not explicitly target violations of semantic protocol requirements derived from specifications.

Despite these advances, fuzzing-based techniques generally offer no guarantees about the absence of protocol-level violations beyond the detection of runtime errors such as crashes, assertion failures, or hangs.
Even when stateful fuzzing reaches deep protocol states, violations that do not manifest as immediate runtime failures may remain undetected.

Protocol state fuzzing based on model learning represents a related line of work.
By inferring abstract state machine models, often in the form of Mealy machines, such techniques can uncover logical flaws in control flow and message ordering~\cite{testing-via-learning,Vaandrager:cacm}.
This approach has been successfully applied to a variety of protocols, including DTLS~\cite{DTLS@USENIX-20, DTLS-Fuzzer@CSCN-25}, TCP~\cite{FH2017,FJV2016}, SSH~\cite{SSH@SPIN-2017}, OpenVPN~\cite{OpenVPN@EuroSPW-2018}, IPSec~\cite{IPSec@IEEEAccess-2019}, QUIC~\cite{Prognosis@SIGCOMM-21,rasool2019}, and EDHOC~\cite{EDHOC-Fuzzer@ISSTA-23}.
These techniques are effective at exposing inconsistencies in protocol state machines, but they typically operate in a black-box manner and do not directly reason about individual protocol requirements or packet-level field constraints.

\section{Symbolic Execution for Protocol Implementations}
Symbolic execution (SE) was introduced in the 1970s as a program analysis technique that explores program behavior by treating selected inputs as symbolic rather than concrete values~\cite{King@CACM-76}.
Advances in constraint solving and tool support over the past two decades have turned SE into a practical and powerful testing technique, capable of uncovering deep bugs and security vulnerabilities that are difficult to expose using random or coverage-guided testing alone~\cite{Cadar:3decades,KLEE@OSDI-08,S2E@12}.

Applying SE to network protocol implementations poses challenges that are not present in traditional single-input programs.
Protocol implementations are often stateful, process sequences of messages rather than isolated inputs, and interact with other protocol parties and the environment.
As a result, SE must account for both the values of packet fields and the order in which packets are exchanged.

Several approaches apply SE to protocol implementations by designating specific fields in incoming packets as symbolic, thereby enabling exploration of code paths reachable under different packet contents~\cite{KleeNet,KleeNet-poster,PedrosaFKGMM15,RathSW18,SymbexNet,tempel2023specification, wilson2024analysing}.
In these approaches, SE is typically used to uncover generic runtime errors, such as crashes or memory safety violations, or to detect interoperability issues between protocol parties.
For example, Pedrosa \etal~\cite{PedrosaFKGMM15} use SE to characterize messages that can be sent by one protocol party but rejected as non-compliant by another, while Rath \etal~\cite{RathSW18} apply SE to explore interoperability scenarios in QUIC implementations.

A recurring difficulty in these settings is protocol statefulness.
To reach deep states, an implementation must process a specific sequence of packets, and the correctness of processing a given packet may depend on information accumulated from earlier messages.
To address this, some approaches incorporate explicit models of protocol behavior.
Wen \etal~\cite{wen2017model} employ model learning to infer a protocol state machine, which is then provided as additional input to SE, while Tempel \etal~\cite{tempel2023specification} show that taking protocol state into account can significantly increase test coverage.
These techniques improve path exploration but still focus primarily on coverage and runtime error detection rather than checking compliance with protocol specifications.

Other lines of work symbolically execute multiple protocol parties jointly in order to explore interactions under different synchronization and failure scenarios.
KleeNet~\cite{KleeNet,KleeNet-poster} extends SE to distributed settings by modeling packet loss and node failures using symbolic variables, allowing the exploration of diverse execution scenarios.
Similar ideas appear in SymNet~\cite{Symnet}, which applies SE to integration testing of protocol implementations.
These approaches are effective at identifying interaction and interoperability issues that arise during joint execution of protocol parties, but have limited ability to expose flaws caused by adversarial or malformed inputs, since inputs are typically generated by protocol-compliant peers.

More recent work combines SE and specification-based testing by checking protocol requirements derived from protocol specification documents.
SymbexNet~\cite{SymbexNet} first applies SE to generate concrete test inputs that explore a wide range of code paths, and then replays these inputs on an unmodified implementation while monitoring the resulting input-output behavior against specification-derived rules.
While this separation allows requirements to be checked without modifying the system under test, SE in this setting may still miss requirement-violating inputs when violating and non-violating inputs exercise the same code path.

Several approaches address this limitation by guiding SE explicitly toward requirement-relevant behaviors.
One line of work, exemplified by Paper~\URoman{1}, embeds assertions and assumptions directly into the SUT's code to constrain SE toward inputs that may violate protocol requirements, enabling systematic exploration of requirement-specific corner cases.
Subsequent work, represented by Paper~\URoman{2}, externalizes this requirement-checking logic into monitors that observe packet exchanges and interact with the SE engine, reducing the need for intrusive modifications to the implementation.
These approaches demonstrate that SE can be effectively adapted from a coverage-oriented testing technique into a tool for checking semantic protocol requirements.

Overall, SE offers strong capabilities for systematically exploring protocol implementations under a wide range of inputs.
However, its effectiveness depends on how statefulness, nondeterminism, and correctness criteria are incorporated into the analysis.
Unguided SE tends to emphasize code coverage and runtime errors, while exposing semantic protocol violations often requires more targeted approaches, including explicit consideration of protocol requirements or alternative notions of correctness.

An alternative to specification-based correctness checking is differential testing, which detects discrepancies by comparing the behavior of multiple implementations under the same inputs.
Differential testing has been widely used to uncover semantic bugs in protocol implementations without requiring a complete specification oracle.
Differential symbolic execution, introduced by Person \etal~\cite{DSE@08}, extends this idea by enabling systematic exploration of inputs that lead to behavioral differences across executions.

Several works have applied differential SE by comparing the behavior of multiple implementations rather than checking conformance against a single specification.
Chau \etal~\cite{SymCerts,PKCS} use SE to derive path constraints for certificate validation logic and compare them across implementations to identify discrepancies.
RFCcert~\cite{RFCcert@ICSE-18} combines specification-derived rules with SE to generate test cases that expose differences between X.509 validators.
Other approaches apply differential analysis to restricted protocol components or employ offline comparison of execution models, such as ParDiff~\cite{ParDiff@24} and SymTCP~\cite{SymTCP@NDSS-20}.
These techniques demonstrate the effectiveness of differential reasoning, but typically operate on isolated components or abstracted models rather than complete, stateful protocol interactions.

Together, these lines of work illustrate how advances in SE have progressively expanded the scope of protocol testing, from coverage-oriented exploration to requirement-aware and differential analysis of stateful protocol implementations.


	\chapter{Conclusion}\label{chapter:conclusion}
Testing network protocol implementations for conformance with their specifications remains challenging due to protocol statefulness, complex dependencies across packet sequences, and large input spaces that are difficult to explore exhaustively.
Even minor deviations from specifications can lead to interoperability issues or security vulnerabilities, yet such deviations often escape detection by conventional testing techniques that rely on concrete test cases or incomplete oracles.

This thesis investigated how symbolic execution can be adapted to systematically and practically address the challenges of testing stateful protocol implementations.
It introduced a requirement-driven approach that uses symbolic execution to explore behaviors relevant to specification requirements, enabling the detection of violations that depend on specific corner-case values of protocol fields.
Building on this foundation, the thesis presented a monitor-based testing technique that externalizes requirement checking from the system under test, allowing protocol requirements to be expressed independently of specific implementations and reused across implementations and versions.
Finally, the thesis explored differential symbolic execution as a complementary approach that removes the need for a fully specified oracle by comparing observable behaviors across implementations, thereby enabling the detection of discrepancies that may indicate nonconformance issues even when explicit requirements have not been encoded.
As a comparative baseline, the thesis also empirically studied the use of coverage-guided fuzzing for testing protocol requirements and contrasted its effectiveness with symbolic execution in the context of CoAP implementations, highlighting the complementary strengths and limitations of the two techniques.

The proposed techniques were evaluated on multiple real-world implementations of the DTLS, CoAP, and QUIC protocols.
Across these studies, the techniques uncovered \emph{72} previously unknown specification violations.
These findings were validated in practice, with the majority of the reported issues fixed by developers.
The discovery of previously unknown bugs in widely deployed protocol implementations indicates that the proposed techniques are applicable beyond controlled academic settings and can be effectively used to analyze production-quality software.
Overall, the results demonstrate that symbolic execution, when guided by protocol requirements and combined with monitor-based and differential analysis, can systematically expose semantic deviations from protocol specifications that are difficult to detect using conventional testing techniques.
	\chapter{Future Work}\label{chapter:future}
	\chapter{Contributions}\label{chapter:contributions}
	\selectlanguage{swedish}

\chapter{Sammanfattning På Svenska}\label{chapter:sammanfattning}
Internet spelar en central roll i dagens samhälle, och möjliggör att vi kan kommunicera och åtnjuta viktiga tjänster som t.ex.\ surfning, bankärenden och social medier. Internet bygger på ett antal nätverksprotokoll, som är implementerade i enheter anslutna till internet. Sådana implementeringar kan innehålla buggar, bland vilka en del, så kallade sårbarheter, kan utnyttjas för att orsaka skada. Ett exempel är den nu ökända Heartbleed-sårbarheten i en vida spridd implementering av TLS-protokollet, some ledde till att flera miljoner patienters data blottades. För att förebygga och åtgärda sådana sårbarheter har tekniker utvecklats för att säkra implementeringar av nätverksprotokoll, varav den dominerande är testning.

Att testa implementationer av nätverksprotokoll mot sina specifikationer är en utmanande uppgift på grund av protokollens tillståndsberoende beteende, komplexa beroenden mellan paketsekvenser och stora indatarymder.
Även små avvikelser från en specifikation kan leda till interoperabilitetsproblem eller allvarliga säkerhetssårbarheter, men sådana fel upptäcks ofta inte med traditionella testmetoder som bygger på konkreta testfall eller ofullständiga testorakel.
Symbolisk exekvering är en testningsteknik som addresserar problemet med stora indatarymder. Den analyserar program i vilka (några av) variabler är \emph{symboliska},
utforskar de kodvägar som är möjliga för specifika värden av de symboliska variablerna, och genererar konkreta för vilka exekveringen följer specifika kodvägar.
Denna avhandling undersöker hur symbolisk exekvering kan anpassas för att systematiskt och praktiskt testa implementationer av tillståndsberoende nätverksprotokoll. 

Avhandlingens första bidrag är en kravdriven metod baserad på symbolisk exekvering, i vilken krav hämtas från protokollets specifikation och uttrycks som  egenskaper hos paketsekvenser, samt realiseras som assertions inlagda i den testade implementeringens källkod. Testningen förbereds genom att spela in och spara en sekvens indatapaket i en interaktion mellan den testade implementeringen och dess motpart. Denna sekvens av indatapaket används sedan som indata under den symboliska exekveringen.
Den symboliska exekveringen styrs mot exekveringsvägar som är relevanta för att bryta mot dessa krav. Genom att selektivt göra delar av indatapaket symboliska och begränsa exekveringen till kravrelevanta scenarier möjliggörs upptäckt av subtila fel som endast uppstår för specifika variabelvärden eller hörnfall.

Det andra bidraget är en monitorbaserad testteknik i vilken implementering och kontroll av krav separeras från det system som testas. Detta görs med hjälp av monitorer: dessa är komponenter testomgivningen som är separerade från den testade implementeringen.
Protokollkrav specificeras i monitorer; dessa observerar paketutbyten, upprätthåller tillståndsinformation, styr den symboliska exekveringen, och kontrollerar krav vid tillämpliga punkter i exekveringen, utan att kräva genomgripande modifieringar av implementeringens källkod.
Genom att protokollkraven inte sätts in i den testade implementeringens källkrav (som i det första bidraget) utan i återanvändbara monitorer,
så förbättras modularitet, möjliggörs återanvändning av krav mellan olika implementationer och versioner, samt förenklas utvidgning till nya protokollkrav.

Det tredje bidraget utforskar differentiell symbolisk exekvering som en testmetod som inte använder explicit formulerade krav.. 
I stället för att kontrollera efterlevnad mot en formellt specificerad kravuppsättning jämförs det observerbara beteendet mellan olika implementationer eller versioner av samma protokoll under identiska symboliska indatasekvenser. Differentiell testning är en etablierad teknik inom vanlig (icke-symbolisk) testning. I bidraget visar vi hur denna teknik kan anpassas för symbolisk exekvering av kommunikationsprotokoll. Anpassningen bygger till vissa delar på den monitorbaserade tekniken som utvecklats i det andra bidraget. Avvikelser i beteende kan därigenom avslöja potentiella fel eller specifikationsbrott även när explicita krav inte har formulerats.

I det fjärde bidraget studeras kravbaserad testning av IoT-protokoll genom en fallstudie utförd på IoT-protokollet CoAP.
Studien är en empirisk jämförelse mellan täckningsstyrd fuzzing och symbolisk exekvering, uförd på CoAP. Den bygger på den kravdrivna testmetoden som utvecklats i avhandlingens första bidrag. Bidraget visar hur denna metod även kan användas för att testa mot protokollkrav genom täckningsstyrd fuzzing.
Analysen belyser teknikernas kompletterande styrkor och svagheter vid testning av protokollkrav.

De föreslagna teknikerna har utvärderats på verkliga implementationer av de vitt använda protokollen DTLS, QUIC och CoAP. 
Some resultat av studierna har ett stort antal tidigare okända buggar och brott mot specifikationskrav upptäckts, av vilka majoriteten har bekräftats och åtgärdats av utvecklarna.
Vidare har arbetet reulterat i en allmänt tillgänglig infrastruktur med återanvändbara testomgivningar och monitorer för de testade protokollen, som kan användas för att reproducera våra resultat och bygga vidare på dem.
Resultaten visar att kravdriven, monitorbaserad och differentiell symbolisk exekvering kan avslöja semantiska avvikelser från protokollspecifikationer som är svåra att upptäcka med konventionella testmetoder, och att teknikerna är praktiskt tillämpbara även på produktionsmogen programvara. 

\section{Finansiering}
Dessa arbeten är delvis finansierade av Stiftelsen för Strategisk Forskning (SSF) genom projektet \href{https://assist-project.github.io/}{aSSIsT} samt av Vetenskapsrådet.

\selectlanguage{english}

	\chapter{Acknowledgements}

I would like to begin by expressing my sincere gratitude to my main supervisor, Bengt Jonsson.
Throughout my doctoral studies, he was always approachable and willing to help, offering thoughtful guidance and constructive feedback whenever it was needed.
At the same time, he placed great trust in me and gave me the freedom to explore ideas, take risks, and shape my own research direction without micromanagement.
This balance of support, trust, and independence was deeply appreciated and played a crucial role in my development as a researcher.

I would like to warmly thank Konstantinos Sagonas, a mentor, close collaborator, and co-author, for his exceptional attention to detail and his dedication to high-quality research.
He has an impressive ability to spot even the smallest issues, and his careful scrutiny consistently strengthened our work.
At the same time, he combined rigor with kindness, making discussions both demanding and genuinely enjoyable.
I frequently turned to him for his opinion to ensure that nothing important had been overlooked, and his insight and thoroughness had a substantial impact on the quality of this thesis.
Beyond his academic guidance, he was also consistently supportive and helpful with practical, everyday matters, which made the entire process smoother and more manageable.

I am deeply grateful to my second supervisor, Paul Fiter\u{a}u-Bro\c{s}tean, for his constant support throughout my PhD.
From the very beginning, he was always there to help with everything from navigating practical matters, such as banking and administrative issues, to engaging in daily discussions about research ideas and technical details.
Beyond his academic guidance, his generosity with his time, his encouragement, and his friendship made the PhD journey far more enjoyable.

I would also like to thank several senior colleagues at the IT Department for their support and guidance throughout my PhD.
In particular, I am grateful to Mohamed Faouzi Atig, Deputy Head of the IT Department, as well as Philipp Rümmer and Parosh Aziz Abdulla, for their help with structuring my PhD studies and for their advice on navigating departmental responsibilities and long-term planning.
Their support was invaluable and greatly contributed to making my doctoral studies run smoothly.

I would also like to thank my officemates and friends, Magnus and Sarbojit, who were both knowledgeable colleagues and genuinely pleasant companions.
Our conversations, ranging from research discussions to broader reflections on life, made everyday work more enjoyable.
I am also grateful to my colleagues Clément, Fredrik, and Stephan, whose presence and interactions contributed to a friendly and supportive environment and made my PhD experience more pleasant overall.

I would also like to thank my friends in Uppsala who made life during my PhD both richer and more enjoyable.
Their companionship, support, and shared moments provided a welcome balance to academic life.
In particular, I am grateful to, in lexicographic order, Alireza, Ashkan, Behnaz, Ehsan, Jalal, Khaled Kasperi, Lan, Mahdi, Maryam, Meriame, Mojtaba, Parnian, Rashid, Sajjad, Sheida, Sina, Soheil, and Sousan for making my time in Uppsala more pleasant and memorable.

I would also like to thank my fiancée, Sara, for her unwavering support, kindness, and care throughout the final years of my PhD.
We got to know each other during the last two years of my doctoral studies, and during that time she consistently helped me manage stress and find calm during difficult moments.
Her presence was a source of love, warmth, and light, especially during the long and dark Swedish winters.
Her support, understanding, and companionship meant more to me than I can express.

Finally, I would like to express my deepest gratitude to my parents, Hanieh and Mohammad, and to my brother, Armin.
Their constant love, support, and encouragement have been with me throughout this journey and in every stage of my life.
I would not be who I am today without them, and this work would not have been possible without their unwavering belief in me and my path.

Thank you for taking the time to read these paragraphs.
If you were part of my PhD journey and were not mentioned, I apologize for any unintended omission.

\section{Funding}

The research in this thesis was partially funded by the Swedish Foundation for Strategic Research (SSF) through project \href{https://assist-project.github.io/}{aSSIsT} and by the Swedish Research Council (Vetenskapsr{\aa}det).


\backmatter
    % References
    % No restriction is set to the reference styles
    % Save your references in References.bib
    \nocite{*} % Remove this for your own citations
    \bibliographystyle{plain}
    \bibliography{References}

\end{document}